\section{Agent report: finite-cutoff estimates and counterexamples}
\label{app-finite:report}

This appendix gives independent proofs of the cutoff estimates and the
neutral-current counterexample used in the main text.  Its conclusions concern
the specified, fixed canonical-momentum cutoff.  They are not claims about every
possible gauge-covariant regularization of QED.  In particular, the large-volume
counterexample below disproves extensivity in the stated free-vacuum
normalization; it does not, by itself, disprove a bound measured relative to the
interacting vacuum.

\subsection{Conventions and finite-regulator realization}
\label{app-finite:conventions}

Set $\hbar=c=1$, let $V=L^3$, and use Heaviside--Lorentz electromagnetic
normalization, so $e^2=4\pi\alpha$.  The scalar fermion momentum space is
\[
 \mathcal P_{\Lambda,L}
 =\operatorname{span}\{V^{-1/2}e^{ip\cdot x}:
 p\in(2\pi/L)\mathbb Z^3,\ |p|\le\Lambda_f\},
 \qquad n=\dim\mathcal P_{\Lambda,L},\quad d=n/V.
\]
Write $P$ for its projection.  The bare one-particle Dirac space is
$\mathbb C^4\otimes\mathcal P_{\Lambda,L}$, of dimension $4n$.  The Dirac
matrices are denoted by $\mathsf a_i$ and $\mathsf b$, and
\[
 D_0=\mathsf a\cdot(-i\nabla)+m\mathsf b,
 \qquad \Gamma_0=\mathbf1_{(-\infty,0)}(D_0),
 \qquad \operatorname{rank}\Gamma_0=2n.
\]
The finite particle--hole transformation identifies the electron--positron
Fock representation with the bare CAR representation.  Under this
identification,
\[
 H_D=d\Gamma(D_0)-\operatorname{Tr}(D_0\Gamma_0)
     =\sum_{p,s}E(p)(b_{p,s}^\dagger b_{p,s}
                         +d_{p,s}^\dagger d_{p,s})\ge0.
\]
All normal ordering in this appendix is relative to $\Gamma_0$.
For unambiguous signs define the fermionic electric charge density in units of
$e$, and the number current, respectively, by
\[
 \rho(x)=-:\!\psi^\dagger(x)\psi(x)\!:,
 \qquad j_i(x)=:\!\psi^\dagger(x)\mathsf a_i\psi(x)\!:.
\]
Thus $\int\rho=N_p-N_e$ and the transverse interaction is
$H_I=e\int j\cdot A$.  With $Z_{\rm tot}=\sum_i Z_i$, the neutral sector obeys
\begin{equation}
 N_e-N_p=Z_{\rm tot},\qquad 0\le Z_{\rm tot}\le2n.
 \label{app-finite:capacity}
\end{equation}
In particular, a nonintegral $Z_{\rm tot}$ gives no neutral sector in this
representation.  Physical integral nuclear charges satisfy the integrality
condition automatically, but still have to satisfy the capacity constraint.

Take transverse photon momenta $0<|k|\le\Lambda_\gamma$, with the constant
photon mode omitted.  The omitted Coulomb zero mode is fixed by the neutral
Gauss constraint, not by inserting a compensating background.  Use the
zero-mean periodic Green function
\[
 G_L(x)=\frac1V\sum_{k\ne0}\frac{4\pi}{|k|^2}e^{ik\cdot x}.
\]
The following nuclear self-energy convention is used throughout this
appendix: remove the entire diagonal contribution of each individual nucleus
and retain $\alpha\sum_{i<j}Z_iZ_jG_L(R_i-R_j)$.  Equivalently, mollify the
nuclei, subtract each mollified nucleus's own periodic Coulomb energy, and
take the point limit.  This removes a finite periodic self term as well as the
divergence.  A convention retaining the periodic Robin self term defines a
different finite-volume energy and must be stated separately.  Distinct
nuclear positions are assumed; collisions have infinite positive mutual
energy under the convention here.

First use the density-product Coulomb operator
\begin{equation}
 H_C^{\rm prod}
 =\frac\alpha2D_G(\rho,\rho)
   +\alpha D_G(\rho,\nu)
   +\alpha\sum_{i<j}Z_iZ_jG_L(R_i-R_j),
 \qquad \nu=\sum_iZ_i\delta_{R_i},
 \label{app-finite:coulomb-definition}
\end{equation}
where $D_G$ is the periodic Coulomb quadratic form with zero mode excluded.
The fermionic Fourier support is contained in $|k|\le2\Lambda_f$, so every
fermionic term in this formula is a bounded finite-dimensional operator.
The nuclei remain point sources.  Their coupling to a band-limited fermion
field does, however, sample only the corresponding band of their Coulomb
potential.  Normal-ordering variants are addressed below.

At fixed $L,\Lambda_f,\Lambda_\gamma$ there are finitely many photon
oscillators of strictly positive frequency.  Every coefficient of their linear
coupling is a bounded fermion matrix.  The creation and annihilation bounds
for these oscillators make $H_I$ infinitesimally $H_\gamma$-bounded.
Consequently Kato--Rellich gives self-adjointness on $D(H_\gamma)$ and
semiboundedness.  The finite-dimensional fermion terms do not change this
domain.  The realization has compact resolvent: $(H_\gamma+1)^{-1}$ is
compact for finitely many positive-frequency oscillators, and the linear
oscillator perturbation is relatively compact.  A ground state therefore
exists in every nonempty charge sector.  These statements are about fixed
regulators and contain no uniform volume estimate.

\subsection{Local exclusion bounds and the transverse-field estimate}
\label{app-finite:photon}

At a fixed point the cutoff CAR are
\[
 \{\psi_a(x),\psi_b^\dagger(x)\}=d\,\delta_{ab},
 \qquad \operatorname{tr}_{\mathbb C^4}\Gamma_0(x,x)=2d.
\]
After dividing by $\sqrt d$, these are four canonical fermion modes.  Hence
$\psi^\dagger\psi$ has spectrum in $[0,4d]$.  Each $\mathsf a_i$ has two
eigenvalues $+1$ and two eigenvalues $-1$.  Inversion symmetry of the momentum
cutoff makes the free-sea current vanish.  Diagonalizing $\mathsf a_i$ in
spin space therefore proves the operator inequalities
\begin{equation}
 -2d\le\rho(x)\le2d,\qquad
 -2d\le j_i(x)\le2d,\qquad
 \sum_{i=1}^3j_i(x)^2\le12d^2.
 \label{app-finite:local-car}
\end{equation}
These bounds already include the pair-creation and pair-annihilation terms.

For a mode with frequency $\omega=|k|$, let $J$ be the relevant transverse
Fourier component of the current and $g=e/(2V\omega)^{1/2}$.  Since photon
operators commute with fermion operators, the exact identity is
\begin{equation}
 \omega a^\dagger a+g(a^\dagger J+aJ^\dagger)
 =\omega\left(a^\dagger+\frac g\omega J^\dagger\right)
             \left(a+\frac g\omega J\right)
       -\frac{g^2}{\omega}J^\dagger J.
 \label{app-finite:square}
\end{equation}
The first term is positive.  This identity does not assert that the shifted
operators for different modes obey canonical commutation relations, and does
not require different currents to commute.  Summing the positive terms,
using the transverse polarization projection, and applying Parseval gives
\begin{align}
 H_\gamma+H_I
 &\ge-\frac{e^2}{2V k_{\min}^2}
       \sum_{k,i}j_i(k)^\dagger j_i(k)\\
 &=-\frac{e^2}{2 k_{\min}^2}
       \int\sum_i j_i(x)^2\,dx
 \ge-\frac{6e^2d^2V}{k_{\min}^2}.
 \label{app-finite:photon-bound}
\end{align}
In the intermediate inequality one may enlarge the sum to all Fourier modes;
this only enlarges the nonnegative quantity being subtracted.  For ordinary
torus photons one may take $k_{\min}=2\pi/L$.  If the model instead has a fixed
infrared exclusion $|k|\ge\kappa>0$, one may take $k_{\min}=\kappa$.  A large
photon ultraviolet cutoff causes no additional loss here because the cutoff
fermion current has Fourier support only up to $2\Lambda_f$.

\subsection{A uniform point-nucleus Coulomb estimate}
\label{app-finite:coulomb}

Let $p_t$ denote the normalized torus heat kernel, set
$G_s=e^{s\Delta}G_L$, $h_s=G_L-G_s$, and
$\nu_s=e^{s\Delta}\nu$.  The heat representation yields
\[
 h_s(x)=4\pi\int_0^s\bigl(p_t(x)-V^{-1}\bigr)\,dt,
 \quad h_s(x)\ge-\frac{4\pi s}{V},
 \quad \|h_s\|_{L^1}\le8\pi s.
\]
The last bound follows from $\|p_t-V^{-1}\|_1\le2$.  Also
\begin{equation}
 0\le G_{2s}(0)
 =\frac1V\sum_{k\ne0}\frac{4\pi e^{-2s|k|^2}}{|k|^2}
 \le C s^{-1/2}.
 \label{app-finite:heat-self}
\end{equation}
For example, the last estimate follows by partitioning the integer lattice
into spherical shells and comparing the resulting Gaussian sum with its
radial integral; its constant is independent of $L$ and $s$.

Expanding the square with the smoothed nuclear density gives the exact
identity
\begin{align}
 H_C^{\rm prod}
 ={}&\frac\alpha2D_G(\rho+\nu_s,\rho+\nu_s)
       +\alpha D_G(\rho,\nu-\nu_s)\nonumber\\
 &+\alpha\sum_{i<j}Z_iZ_jh_{2s}(R_i-R_j)
       -\frac\alpha2G_{2s}(0)\sum_iZ_i^2.
 \label{app-finite:heat-identity}
\end{align}
The first term is positive, including when its operator-valued Fourier
coefficients do not commute, because it is a positive weighted sum of
$T_k^\dagger T_k$.  The other terms satisfy
\begin{align*}
 \alpha D_G(\rho,\nu-\nu_s)
   &\ge-16\pi\alpha d s Z_{\rm tot},\\
 \alpha\sum_{i<j}Z_iZ_jh_{2s}(R_i-R_j)
   &\ge-\frac{4\pi\alpha s}{V}Z_{\rm tot}^2
    \ge-8\pi\alpha d s Z_{\rm tot},\\
 -\frac\alpha2G_{2s}(0)\sum_iZ_i^2
   &\ge-C\alpha s^{-1/2}Z_{\max}Z_{\rm tot}.
\end{align*}
The second line uses the finite-mode neutrality constraint
\eqref{app-finite:capacity}; without it, this particular argument would not
give the displayed estimate.  Taking, for instance, $s=\Lambda_f^{-2}$ proves
\begin{equation}
 H_C^{\rm prod}\ge-C_{\alpha,\Lambda_f,L_0,Z_{\max}}K
 \quad\hbox{for }L\ge L_0>0,
 \label{app-finite:coulomb-bound}
\end{equation}
since $d\le C(\Lambda_f^3+L_0^{-3})$.  No nuclear minimum-separation
assumption is used.

\subsection{Normal-ordering audit and the resulting lower bounds}
\label{app-finite:wick}

Normal ordering the density alone is not the same as fully Wick ordering its
product.  This difference must be specified rather than absorbed without
explanation.  Here it is an extensive, bounded correction.

For a one-particle matrix $A$, write
$X_A=d\Gamma(A)-\operatorname{Tr}(\Gamma_0A)$.  Direct CAR contraction gives
\begin{equation}
 X_A X_B
 =:\!X_A X_B\!:
  +X_{A(1-\Gamma_0)B-B\Gamma_0A}
  +\operatorname{Tr}\bigl(\Gamma_0A(1-\Gamma_0)B\bigr).
 \label{app-finite:wick-identity}
\end{equation}
Here the first term means full Wick ordering relative to the same free sea.
For the density Fourier matrices $A=P e^{ik\cdot x}P$, $B=A^\dagger$, both
matrix norms are at most one.  The middle matrix has norm at most two, while
the one-particle dimension is $4n$.  The estimates
\[
 \|X_M\|\le\|M\|_1+|\operatorname{Tr}(\Gamma_0M)|,
 \qquad
 |\operatorname{Tr}(\Gamma_0A(1-\Gamma_0)B)|\le2n
\]
therefore bound the difference in \eqref{app-finite:wick-identity} by a
constant times $n$, uniformly in $k$.  Finally,
\begin{equation}
 \frac1V\sum_{0<|k|\le2\Lambda_f}\frac1{|k|^2}
 \le C\Lambda_f
 \label{app-finite:coulomb-sum}
\end{equation}
follows from $\sum_{0<|\ell|\le M}|\ell|^{-2}\le CM$ in three dimensions.
It follows that replacing the first term in
\eqref{app-finite:coulomb-definition} by its fully Wick-ordered version
changes the Hamiltonian by an operator of norm at most $C\alpha n\Lambda_f$.
Subtracting just its free-vacuum expectation obeys the same bound.  Both costs
are $O(V)$ for $L\ge L_0$.

Combining \eqref{app-finite:photon-bound},
\eqref{app-finite:coulomb-bound}, and $H_D\ge0$ proves
\begin{equation}
 H\ge-C_{\Lambda,L_0}\bigl[K+V(1+L^2)\bigr].
 \label{app-finite:polynomial-bound}
\end{equation}
With a fixed photon infrared exclusion $|k|\ge\kappa>0$, the stronger
conclusion is
\begin{equation}
 H\ge-C_{\Lambda,L_0,\kappa}(K+V).
 \label{app-finite:ir-extensive}
\end{equation}
Both bounds hold at arbitrary fixed positive coupling, including the physical
value of $\alpha$, with constants depending on that coupling.  A finite fixed
mass counterterm is a bounded fermion operator of norm $O(n)$ and does not
alter these conclusions.  A fixed vacuum-energy-density counterterm is
$O(V)$.  Charge renormalization must be implemented consistently in Coulomb
and transverse couplings; a fixed nonzero renormalized coupling changes the
constants.  No assertion is made here about tuning these counterterms to a
cutoff-independent physical theory.

\subsection{A neutral-current trial state with energy of order minus \texorpdfstring{$L^5$}{L to the fifth power}}
\label{app-finite:current-counterexample}

The $VL^2$ loss in \eqref{app-finite:polynomial-bound} reflects an actual
instability of extensivity in this canonical-momentum cutoff model.  The
following construction is in the exact neutral sector, with no nuclei.

For $q=2\pi/L$ define scalar matrices on $\mathcal P_{\Lambda,L}$ by
\[
 U=P e^{iqx_1}P,
 \qquad C=\frac{U+U^\dagger}{2}=P\cos(qx_1)P.
\]
Choose a self-adjoint unitary $S$ that equals $\operatorname{sign}(C)$ on
$(\ker C)^\perp$, choosing either sign on each zero eigenspace.  The matrix
\begin{equation}
 \gamma=\frac{1+\mathsf a_2S}{2}
 \label{app-finite:trial-projector}
\end{equation}
is an orthogonal projector of rank $2n$.  Its Slater determinant is therefore
an exact charge-zero state, not a grand-canonical mixture.  Its density is
\[
 \operatorname{tr}_{\mathbb C^4}\gamma(x,x)
 =2P(x,x)
 =\operatorname{tr}_{\mathbb C^4}\Gamma_0(x,x).
\]
Thus $\langle\rho(x)\rangle=0$ identically.  Nevertheless its transverse
current has a macroscopic Fourier component:
\begin{align}
 J&=\int\cos(qx_1)\langle j_2(x)\rangle\,dx
   =\operatorname{Tr}\bigl(\mathsf a_2C(\gamma-\Gamma_0)\bigr)
   =2\operatorname{Tr}|C|\nonumber\\
  &\ge2\operatorname{Tr}C^2
   =\#\{p:p\hbox{ and }p+qe_1\hbox{ belong to the cutoff band}\}
   \ge c_{\Lambda_f}V
 \label{app-finite:macroscopic-current}
\end{align}
for all sufficiently large $L$.  The free-sea term in this trace vanishes
because $\Gamma_0$ is diagonal in momentum and $C$ has zero momentum diagonal.
The inequality uses $\|C\|\le1$.  For the final estimate it suffices to count
momenta with $|p|\le\Lambda_f/2$ once $q\le\Lambda_f/2$.

The free Dirac energy of this Slater determinant is bounded above by
$4n\sqrt{m^2+\Lambda_f^2}$.  For every nonzero momentum, set
$U_k=P e^{ik\cdot x}P$, with the spin identity implicit.  The Slater covariance
formula gives
\begin{equation}
 \langle\rho_k^\dagger\rho_k\rangle
 =\operatorname{Tr}\bigl(\gamma U_k^\dagger(1-\gamma)U_k\bigr)
 \le2n,
 \label{app-finite:trial-charge-fluctuation}
\end{equation}
because the mean density is zero.  Equations
\eqref{app-finite:coulomb-sum} and
\eqref{app-finite:trial-charge-fluctuation} bound its density-product Coulomb
energy above by $C\alpha n\Lambda_f$.  The other normal-ordering conventions
in Section~\ref{app-finite:wick} change this by at most another $O(V)$ term.

Now take the tensor product with a photon coherent state whose mean transverse
potential is
\[
 A_2(x)=-a\cos(qx_1),\qquad A_1=A_3=0,
\]
and whose mean transverse electric field vanishes.  This is an allowed
coherent state whenever $q\le\Lambda_\gamma$.  The normally ordered free
photon energy equals its classical magnetic energy,
$Vq^2a^2/4$.  The current coupling is $-eaJ$.  Minimizing their sum over the
real number $a$ yields
\begin{equation}
 \inf_a\left(\frac{Vq^2a^2}{4}-eaJ\right)
 =-\frac{e^2J^2}{Vq^2}.
 \label{app-finite:coherent-minimum}
\end{equation}
Every trial state at finite $L$ has finite photon number and lies in the
Hamiltonian form domain.  Its photon number need not be uniformly bounded as
$L$ increases; the original stability question does not impose such a bound.

The variational principle now proves, for all sufficiently large $L$,
\begin{equation}
 E^{\rm vac}_{\Lambda,L}
 \le C_\Lambda V-c_\Lambda e^2VL^2.
 \label{app-finite:l5-instability}
\end{equation}
Thus no inequality $E^{\rm vac}_{\Lambda,L}\ge-C_\Lambda V$ can hold in the
stated free-vacuum normalization.  This is true at every fixed nonzero
coupling, however small.  The trial state is charge neutral, its mean charge
density vanishes, and its Coulomb fluctuations and pair rest-energy cost are
only extensive.  The negative contribution comes from the coherent
long-wavelength transverse field.

The mechanism depends on the fixed canonical-momentum cutoff: large vector
potentials do not move the retained fermion modes in a gauge-covariant way.
It is therefore inappropriate to extrapolate
\eqref{app-finite:l5-instability} to a different, gauge-covariant regulator
without a new argument.  A vacuum counterterm proportional to $V$ or a fixed
mass counterterm cannot cancel the $L^5$ term.  Subtracting the actual
interacting vacuum energy removes this particular obstruction by definition,
but leaves the nuclear vacuum-relative inequality unproved.  The fixed-volume
positivity-after-a-vacuum-shift statement in \cite{JW2018} is compatible with
this conclusion and supplies no contrary extensive estimate.

\subsection{A separate small-box obstruction to an all-volume assertion}
\label{app-finite:small-box}

Under the nuclear self convention in
Section~\ref{app-finite:conventions}, even a vacuum-relative estimate uniform
over every $L>0$ is false.  Take $L$ so small that
$2\pi/L>\max\{\Lambda_f,\Lambda_\gamma\}$.  Only the fermion momentum $p=0$
remains, and there are no photon modes.  Put two unit nuclei in the box.
Neutrality forces occupation of both zero-momentum electron spin states and
no positrons.  Every nonzero Fourier component of the fermion density
vanishes as an operator.  Hence the electron--nucleus and fermion--fermion
Coulomb terms vanish, including the ordering corrections discussed above,
and
\begin{equation}
 E_{\Lambda,L}(R_1,R_2)-E^{\rm vac}_{\Lambda,L}
 =2m+\alpha G_L(R_1-R_2),\qquad E^{\rm vac}_{\Lambda,L}=0.
 \label{app-finite:small-box-energy}
\end{equation}
The zero-mean Green function $G_1$ has a point $r_0\ne0$ where $G_1(r_0)<0$:
it is integrable, is positive near its Coulomb singularity, and has integral
zero.  Choose $R_1-R_2=Lr_0$ on the torus.  Scaling gives
$G_L(Lr_0)=L^{-1}G_1(r_0)$, so the right-hand side of
\eqref{app-finite:small-box-energy} tends to $-\infty$ as $L\downarrow0$ at
fixed $K=2$.  A fixed mass renormalization only changes the finite $2m$ term,
and a shared vacuum-energy subtraction cancels from the energy difference.

This small-box example uses the specified periodic self-energy convention.
If a finite periodic nuclear self term is retained instead, its contribution
must be included before asserting an analogous formula.  It is not legitimate
to change that convention in the middle of a proof.  A macroscopic theorem
can avoid the small-box issue by explicitly restricting $L\ge L_0>0$.

\subsection{Scope of the conclusions}
\label{app-finite:scope}

The finite-regulator self-adjointness statement,
\eqref{app-finite:polynomial-bound}, the modified-model extensive bound
\eqref{app-finite:ir-extensive}, and both trial-state obstructions are proved
above.  Their proofs do not use a no-pair stability theorem.  Pair creation is
retained throughout through the exact finite particle--hole representation.
The extensive estimate with a fixed photon infrared exclusion is a theorem
about that additional regulator, and is not uniform as $\kappa\downarrow0$.

For $L\ge L_0$, the neutral-current example neither proves nor disproves
$E_{\Lambda,L}(R,Z)-E^{\rm vac}_{\Lambda,L}\ge-C_\Lambda K$.
Controlling this difference would require an estimate comparing the
interacting vacua with and without sources, rather than separate absolute
lower bounds.  None of the estimates here gives that comparison.  Likewise,
they establish no ultraviolet limit, physical mass or charge identification,
thermodynamic limit relative to the interacting vacuum, binding theorem, or
crystallization theorem.
